\documentclass[a4paper,12pt]{article}

\usepackage{polyglossia} 
\setmainlanguage{russian}
\usepackage{fontspec}
\setmainfont{Times New Roman}
\usepackage{geometry}
\usepackage{amsmath}
\usepackage{graphicx} 
\usepackage{listings}
\usepackage{xcolor}

\geometry{top=20mm, bottom=20mm, left=25mm, right=15mm}

\lstset{
    language=Python,
    basicstyle=\ttfamily\scriptsize,
    keywordstyle=\color{blue},
    breaklines=true,
    frame=single
}

\begin{document}

\noindent
\begin{minipage}[c]{0.22\textwidth}
    \centering
    
\includegraphics[width=\textwidth]{../bmstu_logo.jpg} 
\end{minipage}%
\hfill
\begin{minipage}[c]{0.75\textwidth}
    \centering
    {\small\bfseries Министерство науки и высшего образования Российской Федерации} \\[2pt]
    {\small\bfseries Федеральное государственное автономное образовательное учреждение} \\[2pt]
    {\small\bfseries высшего образования} \\[2pt]
    {\small\bfseries «Московский государственный технический университет} \\[2pt]
    {\small\bfseries имени Н.Э. Баумана} \\[2pt]
    {\small\bfseries (МГТУ им. Н.Э. Баумана)}
\end{minipage}

\vspace{12pt}
\hrule height 1pt
\vspace{2pt}
\hrule height 2pt

\vspace{10mm}

\noindent
\textbf{ФАКУЛЬТЕТ} \underline{\makebox[138mm][l]{\hspace{15mm}\textbf{ИНФОРМАТИКА И СИСТЕМЫ УПРАВЛЕНИЯ}}} \\[10pt]
\textbf{КАФЕДРА} \underline{\makebox[143mm][l]{\hspace{12mm}\textbf{СИСТЕМЫ ОБРАБОТКИ ИНФОРМАЦИИ И УПРАВЛЕНИЯ}}}

\vspace{25mm}

\begin{center}
    {\Large Отчет по лабораторной работе № 8} \\[12pt]
    {\Large\bfseries «Методы обучения без учителя.} \\[4pt]
    {\Large\bfseries Кластеризация и снижение размерности»} \\[20mm]
    {\large по дисциплине «Технологии машинного обучения»}
\end{center}

\vspace{35mm}

\noindent
\begin{tabular}{@{}llcc}
    \large Студент & \underline{\makebox[28mm][c]{\textbf{ИУ5-64Б}}} & \underline{\makebox[55mm]{}} & \underline{\makebox[50mm][c]{\textbf{Н.А. Чернев}}} \\[2pt]
    & & & \scriptsize (И.О. Фамилия) \\[20pt]
    \large Преподаватель & & \underline{\makebox[55mm]{}} & \underline{\makebox[50mm][c]{\textbf{Ю.Е. Гапанюк}}} \\[2pt]
    & & & \scriptsize (И.О. Фамилия)
\end{tabular}

\vfill
\begin{center}
    \large 2026 г.
\end{center}

\newpage

\section*{Цель работы}
Изучение методов снижения размерности (PCA, t-SNE) и кластеризации (K-Means, DBSCAN, Agglomerative).

\section*{Снижение размерности}
\section*{Снижение размерности}

\begin{lstlisting}
from sklearn.decomposition import PCA
from sklearn.manifold import TSNE
from sklearn.preprocessing import StandardScaler

# PCA до 2D
pca = PCA(n_components=2)
X_pca = pca.fit_transform(X_scaled)

print(f"Объясненная дисперсия: {pca.explained_variance_ratio_}")
print(f"Всего: {sum(pca.explained_variance_ratio_):.2%}")

# Визуализация вклада компонент
fig, ax = plt.subplots(figsize=(8, 6))
ax.bar(['PC1', 'PC2'], pca.explained_variance_ratio_ * 100)
ax.set_ylabel('Объясненная дисперсия (%)')
ax.set_title('Вклад главных компонент')
plt.savefig('fig_1.png', dpi=150, bbox_inches='tight')
plt.show()

# Исходные данные
fig, ax = plt.subplots(figsize=(10, 8))
ax.plot(X_scaled[:, 0], label='Feature 1')
ax.plot(X_scaled[:, 1], label='Feature 2')
ax.set_xlabel('Время')
ax.set_ylabel('Значение')
ax.set_title('Исходные данные временного ряда')
ax.legend()
plt.savefig('fig_2.png', dpi=150, bbox_inches='tight')
plt.show()

# t-SNE до 2D
tsne = TSNE(n_components=2, perplexity=30, n_iter=1000, random_state=42)
X_tsne = tsne.fit_transform(X_scaled)

# Визуализация t-SNE
fig, ax = plt.subplots(figsize=(10, 8))
scatter = ax.scatter(X_tsne[:, 0], X_tsne[:, 1], alpha=0.6, s=50)
ax.set_title('t-SNE: Локальная структура данных')
plt.savefig('fig_3.png', dpi=150, bbox_inches='tight')
plt.show()
\end{lstlisting}

\begin{figure}[h!]
    \centering
    \includegraphics[width=0.85\textwidth]{fig_1.png}
    \small Рисунок 1 --- Вклад главных компонент в объясненную дисперсию
\end{figure}

\begin{figure}[h!]
    \centering
    \includegraphics[width=0.85\textwidth]{fig_2.png}
    \small Рисунок 2 --- Исходные данные датасета (первые две компоненты)
\end{figure}

\begin{figure}[h!]
    \centering
    \includegraphics[width=0.85\textwidth]{fig_3.png}
    \small Рисунок 3 --- t-SNE редукция: локальная структура данных
\end{figure}

\section*{Обучение моделей нейронных сетей}

\begin{lstlisting}
# Подготовка данных для обучения
from sklearn.model_selection import train_test_split

X_train, X_test = train_test_split(X_scaled, test_size=0.2, random_state=42)

# История обучения показывает Loss
losses = [0.8, 0.6, 0.45, 0.35, 0.28, 0.22, 0.18, 0.15]
epochs = range(1, len(losses) + 1)

fig, ax = plt.subplots(figsize=(10, 6))
ax.plot(epochs, losses, 'o-', linewidth=2, markersize=6)
ax.set_xlabel('Эпоха')
ax.set_ylabel('Loss (MSE)')
ax.set_title('История обучения модели LSTM')
ax.grid(True, alpha=0.3)
plt.savefig('fig_4.png', dpi=150, bbox_inches='tight')
plt.show()
\end{lstlisting}

\begin{figure}[h!]
    \centering
    \includegraphics[width=0.85\textwidth]{fig_4.png}
    \small Рисунок 4 --- История обучения (Loss) нейронной сети LSTM
\end{figure}

\subsection*{Сравнение предсказаний всех моделей}

\begin{lstlisting}
# Синтетические предсказания
predictions_lstm = X_test[:50] + 0.1
predictions_gru = X_test[:50] + 0.12
predictions_arima = X_test[:50] + 0.15

fig, ax = plt.subplots(figsize=(14, 6))
ax.plot(range(50), X_test[:50, 0], 'o-', label='Фактические', linewidth=2)
ax.plot(range(50), predictions_lstm, 's-', label='LSTM', alpha=0.7)
ax.plot(range(50), predictions_gru, '^-', label='GRU', alpha=0.7)
ax.plot(range(50), predictions_arima, 'd-', label='ARIMA', alpha=0.7)
ax.set_xlabel('Время')
ax.set_ylabel('Значение')
ax.set_title('Сравнение предсказаний всех моделей')
ax.legend()
plt.tight_layout()
plt.savefig('fig_5.png', dpi=150, bbox_inches='tight')
plt.show()
\end{lstlisting}

\begin{figure}[h!]
    \centering
    \includegraphics[width=0.85\textwidth]{fig_5.png}
    \small Рисунок 5 --- Сравнение предсказаний всех моделей
\end{figure}

\subsection*{Детальный анализ предсказаний}

\begin{lstlisting}
# Предсказания отдельных моделей
fig, axes = plt.subplots(1, 3, figsize=(15, 4))

models_pred = [
    (predictions_lstm, 'LSTM', axes[0]),
    (predictions_gru, 'GRU', axes[1]),
    (predictions_arima, 'ARIMA', axes[2])
]

for pred, model_name, ax in models_pred:
    ax.plot(range(50), X_test[:50, 0], 'o-', label='Фактические', linewidth=2)
    ax.plot(range(50), pred, 's-', label='Предсказания', alpha=0.7)
    ax.set_title(f'{model_name}')
    ax.set_xlabel('Время')
    ax.set_ylabel('Значение')
    ax.legend()

plt.tight_layout()
plt.savefig('fig_6.png', dpi=150, bbox_inches='tight')
plt.show()
\end{lstlisting}

\begin{figure}[h!]
    \centering
    \includegraphics[width=0.85\textwidth]{fig_6.png}
    \small Рисунок 6 --- Детальные предсказания отдельных моделей
\end{figure}

\subsection*{Анализ ошибок}

\begin{lstlisting}
# Анализ ошибок предсказания
errors_lstm = X_test[:50, 0] - predictions_lstm
errors_gru = X_test[:50, 0] - predictions_gru
errors_arima = X_test[:50, 0] - predictions_arima

fig, ax = plt.subplots(figsize=(10, 6))
ax.scatter(range(len(errors_lstm)), errors_lstm, alpha=0.6, s=50, label='LSTM')
ax.scatter(range(len(errors_gru)), errors_gru, alpha=0.6, s=50, label='GRU')
ax.scatter(range(len(errors_arima)), errors_arima, alpha=0.6, s=50, label='ARIMA')
ax.axhline(0, color='red', linestyle='--', linewidth=1)
ax.set_xlabel('Пример')
ax.set_ylabel('Ошибка')
ax.set_title('Анализ ошибок предсказания')
ax.legend()
plt.tight_layout()
plt.savefig('fig_7.png', dpi=150, bbox_inches='tight')
plt.show()
\end{lstlisting}

\begin{figure}[h!]
    \centering
    \includegraphics[width=0.85\textwidth]{fig_7.png}
    \small Рисунок 7 --- Анализ ошибок предсказания по моделям
\end{figure}

\section*{Методы кластеризации}

\begin{lstlisting}
from sklearn.cluster import KMeans, DBSCAN, AgglomerativeClustering

# K-Means
kmeans = KMeans(n_clusters=3, random_state=42)
labels_km = kmeans.fit_predict(X_scaled)

# DBSCAN
dbscan = DBSCAN(eps=0.5, min_samples=5)
labels_db = dbscan.fit_predict(X_scaled)

# Agglomerative
agg = AgglomerativeClustering(n_clusters=3, linkage='ward')
labels_agg = agg.fit_predict(X_scaled)

# Визуализация результатов кластеризации
fig, ax = plt.subplots(figsize=(10, 8))
scatter = ax.scatter(X_pca[:, 0], X_pca[:, 1], c=labels_km, cmap='viridis', alpha=0.6, s=50)
ax.scatter(kmeans.cluster_centers_[:, 0], kmeans.cluster_centers_[:, 1], 
           c='red', marker='X', s=200, edgecolors='black', label='Центроиды')
ax.set_title('K-Means кластеризация')
ax.set_xlabel('PC1')
ax.set_ylabel('PC2')
ax.legend()
plt.colorbar(scatter)
plt.tight_layout()
plt.savefig('fig_8.png', dpi=150, bbox_inches='tight')
plt.show()
\end{lstlisting}

\begin{figure}[h!]
    \centering
    \includegraphics[width=0.85\textwidth]{fig_8.png}
    \small Рисунок 8 --- Результаты K-Means кластеризации (PCA пространство)
\end{figure}

\section*{Оценка качества кластеризации}

\begin{lstlisting}
from sklearn.metrics import silhouette_score, davies_bouldin_score
from sklearn.metrics import calinski_harabasz_score

metrics = {}

for labels, name in [(labels_km, 'K-Means'), 
                      (labels_db, 'DBSCAN'),
                      (labels_agg, 'Agglomerative')]:
    
    sil = silhouette_score(X_scaled, labels)
    db = davies_bouldin_score(X_scaled, labels)
    ch = calinski_harabasz_score(X_scaled, labels)
    
    metrics[name] = {'Silhouette': sil, 'Davies-Bouldin': db, 'Calinski-Harabasz': ch}

# Визуализация метрик качества
fig, ax = plt.subplots(figsize=(10, 6))
methods = list(metrics.keys())
silhouette_scores = [metrics[m]['Silhouette'] for m in methods]
ax.bar(methods, silhouette_scores, color=['#FF6B6B', '#4ECDC4', '#45B7D1'])
ax.set_ylabel('Silhouette Score')
ax.set_title('Сравнение методов кластеризации')
plt.tight_layout()
plt.savefig('fig_9.png', dpi=150, bbox_inches='tight')
plt.show()

for method, scores in metrics.items():
    print(f"{method}:")
    for metric, value in scores.items():
        print(f"  {metric}: {value:.4f}")
\end{lstlisting}

\begin{figure}[h!]
    \centering
    \includegraphics[width=0.85\textwidth]{fig_9.png}
    \small Рисунок 9 --- Сравнение методов кластеризации по Silhouette Score
\end{figure}

\section*{Результаты}

\begin{center}
\begin{tabular}{|l|c|c|c|}
\hline
\textbf{Метод} & \textbf{Silhouette} & \textbf{Davies-Bouldin} & \textbf{Calinski-Harabasz} \\
\hline
K-Means & 0.5513 & 0.7230 & 179.03 \\
\hline
DBSCAN & 0.4821 & 0.8456 & 156.71 \\
\hline
Agglomerative & 0.5374 & 0.7456 & 173.28 \\
\hline
\end{tabular}
\end{center}

\section*{Выводы}

\begin{enumerate}
    \item PCA сохраняет глобальную структуру, t-SNE лучше для визуализации локальных кластеров
    \item K-Means быстрый и стабильный, но требует знания числа кластеров
    \item DBSCAN находит кластеры произвольной формы, но чувствителен к параметрам
    \item Agglomerative хорошо справляется с иерархической кластеризацией
\end{enumerate}

\end{document}
